% III. Methodology (continued): formulation, models and evaluation protocol
% Subsections H to J. Based on the model selection section of Milestone 2
% and on docs/resultados/ (metadata.json of each model).

% ---------------------------------------------------------------------
\subsection{Problem Formulation and Metrics}
\label{sec:formulation}
The task is supervised multiclass classification: given the ROI $x$ of a disc, the model predicts a class $\hat{y}\in\{\text{Degeneration},\ \text{Normal},\ \text{Protrusion}\}$. The classes are nominal: the priority used during labeling (Section~\ref{sec:annotation}) is not a severity order, so the task was not posed as regression or ordinal classification.

The metrics are accuracy, per-class precision, recall and F1 and their macro averages, balanced accuracy (equal to macro recall), macro one-versus-rest area under the ROC curve, and log loss. Accuracy alone is misleading under this imbalance: a classifier that always predicted Protrusion would get 73\% of the original test ROIs right. The main comparison criterion is therefore macro F1, together with per-class recall \cite{mosquera2024}. To answer Q3, the clinical false-negative rate is also reported: the fraction of pathological discs (Degeneration or Protrusion) that the model classifies as Normal.

% ---------------------------------------------------------------------
\subsection{Models and Training}
\label{sec:models}
Four models were chosen that meet four criteria: (i) they can classify an ROI into three classes; (ii) they have precedents in lumbar MRI or in comparable medical tasks; (iii) they represent different families, so that the comparison is meaningful; and (iv) they can be trained with the same patient-level split and evaluated with the same metrics. All of them start from ImageNet-pretrained weights. The 128$\times$128~px grayscale ROIs are converted to three channels, resized to 224$\times$224~px and scaled to $[0,1]$; for ResNet-50 they are also normalized with the ImageNet mean and standard deviation. Table~\ref{tab:hyper} summarizes each model's configuration.

\begin{table*}[t]
\centering\footnotesize
\caption{Configuration of the four models. Values chosen among several options were set by validation macro F1 (or accuracy, for YOLO)}
\label{tab:hyper}
\begin{tabular}{>{\raggedright\arraybackslash}p{0.15\textwidth}>{\raggedright\arraybackslash}p{0.2\textwidth}>{\raggedright\arraybackslash}p{0.19\textwidth}>{\raggedright\arraybackslash}p{0.36\textwidth}}
\toprule
Model & Input & Optimization and loss & Hyperparameters \\
\midrule
1. YOLOv8s-cls & 224$\times$224 image & AdamW; cross-entropy & Initial rate $5\times10^{-4}$, final $\times 0.1$, cosine with 1 warm-up epoch; weight decay 0.01; 4 frozen layers; dropout 0.3; batch 64; up to 15 epochs, patience 4 \\
2. ResNet-50 & 224$\times$224 image, ImageNet normalization & AdamW; cross-entropy with 0.1 label smoothing & Rate $10^{-4}$ with cosine decay; weight decay 0.01; dropout 0.3; batch 64; up to 10 epochs, patience 4; best epoch: 2 \\
3. SVM & 2,048 features from the fine-tuned ResNet-50 $\to$ PCA 256 & Hinge loss (SVC) & RBF kernel, $\gamma$ = \texttt{scale}, one-versus-rest; $C\in\{0.1, 0.5, 2, 8\}$, chosen $C=0.1$ \\
4. Random Forest & Same as the SVM & Gini impurity & 300 trees; maximum depth $\in\{10, 20, \text{none}\}$, chosen 20 \\
\bottomrule
\end{tabular}
\end{table*}

\subsubsection{YOLOv8s-cls}
It is the classification variant of the YOLO family, the most common in the reviewed literature \cite{liawrungrueang2023,prisilla2023,wang2025}, and it is consistent with the rest of the pipeline, which already uses YOLO to detect the discs. It was fine-tuned with AdamW, cross-entropy, the first four layers frozen and a dropout of 0.3. Since the main augmentation is already in the files, online augmentation was limited to varying the crop aspect ratio (from 3:4 to 4:3) and brightness ($\pm 5\%$), without flips.

\subsubsection{ResNet-50}
It is a convolutional network from a different (residual) family, with precedents in lumbar disc classification \cite{tan2024,alkubaisi2022}. The first convolution and the first residual block were frozen, and the final layer was replaced by a dropout of 0.3 and a linear layer with three outputs. It was trained with AdamW, a cosine-decaying learning rate, cross-entropy with label smoothing of 0.1 and mild online augmentation (rotation of $\pm 6^{\circ}$, 4\% shift, zoom from 0.95 to 1.05, and brightness and contrast of $\pm 10\%$). The best epoch was chosen by validation macro F1, with early stopping.

\subsubsection{SVM and Random Forest on Deep Features}
A classical classifier on raw pixels treats each value separately and does not recognize shape. ResNet-50 was therefore used as a feature extractor: without its final layer, it yields 2,048 values per image, which are standardized, reduced to 256 components with principal component analysis (PCA; 80.0\% of the variance) and classified with an RBF-kernel SVM or a 300-tree Random Forest. The technique has precedents on this same corpus \cite{sari2025} and in other medical tasks \cite{khoiwal2024}; the Random Forest only has precedents on tabular data \cite{zhang2023} and is presented as our own extension. The scaler, the PCA and the classifiers are fitted on the training set only.

The initial design used ResNet-50 with frozen ImageNet weights, so that the SVM and the Random Forest would not depend on the fine-tuned network. On validation they performed poorly (macro F1 of 0.64 and 0.60), so the already fine-tuned ResNet-50 was adopted as the extractor (F1 of 0.80 and 0.79). The cost is that models 3 and 4 depend on model 2. The ImageNet version is kept as an ablation (Section~\ref{sec:res_q2}). As course baselines, a logistic regression and a KNN were also trained on the same features; they are not among the four models.

% ---------------------------------------------------------------------
\subsection{Evaluation Protocol}
\label{sec:protocol}
The four models use the same patient-level split (Table~\ref{tab:split}) and the same seed (42). Epochs, hyperparameters and variants were chosen on validation; the test set was evaluated once, with the final configuration. The variants tried on validation (other seeds, learning rates, smoothing and ensembles) are recorded in the repository.

The course guidelines recommend augmenting only the training set. Here, the validation and test sets were also augmented, to obtain the same count per class (Section~\ref{sec:augmentation}). Validation and test copies never enter training, so there is no leakage; even so, to keep near-identical copies from inflating the metrics, \textbf{the main results are computed on the 768 original test ROIs}, with the natural class proportions (42 Degeneration, 164 Normal and 562 Protrusion). The full, balanced test set with copies (5,253 images) is reported separately. With 768 images, the 95\% confidence interval of accuracy is about $\pm 2.2$ points (Wilson interval).
